% Beamer Presentation: The Path to Automated Narrative Processing
\documentclass[aspectratio=169]{beamer}

% Theme and colors - customize as needed
\usetheme{Madrid}
\usecolortheme{default}

% Packages
\usepackage[utf8]{inputenc}
\usepackage[T1]{fontenc}
\usepackage{graphicx}
\usepackage{booktabs}
\usepackage{hyperref}
\usepackage{tikz}

% Bibliography
\usepackage[backend=bibtex,style=authoryear]{biblatex}
\addbibresource{references.bib}

% Metadata
\title{The Path to Automated Narrative Processing}
\subtitle{Narrative and the Quest for Human-like AI}
\author{Your Name}
\institute{Your Institution}
\date{\today}

\begin{document}

% Title slide
\begin{frame}
  \titlepage
\end{frame}

% Outline
\begin{frame}{Outline}
  \tableofcontents
\end{frame}

%=============================================================================
\section{Introduction}
%=============================================================================

\begin{frame}{The Central Question}
  \begin{block}{How does narrative figure into human-like AI?}
    How does narrative figure into the drive towards artificial intelligences that bear comparison to human cognition?
  \end{block}
  
  \pause
  \vspace{1em}
  
  \textbf{Two existing approaches:}
  \begin{itemize}
    \item Narrative computation — analysis of narrative structures
    \item LLM narrative generation — producing narrative text
  \end{itemize}
  
  \pause
  \vspace{1em}
  
  \alert{Neither resembles human narrative processing.}
\end{frame}

\begin{frame}{Presentation Overview}
  This presentation will:
  \begin{enumerate}
    \item Describe the function of ongoing narrative processing in human intelligence
    \item Consider the state of AI in approximating that functionality
    \item Offer a research roadmap toward human-like event processing
    \item Consider ethical implications and limitations
  \end{enumerate}
\end{frame}

%=============================================================================
\section{The Problem Narrative Processing Solves}
%=============================================================================

\begin{frame}{Orienting, Framing, Parsing}
  Ongoing narrative processing helps us:
  
  \begin{itemize}
    \item<1-> \textbf{Orient} ourselves towards experience
    \item<2-> \textbf{Frame} responses to experience
    \item<3-> \textbf{Retroactively parse} experience
  \end{itemize}
  
  \vspace{1em}
  \only<4>{
    \begin{block}{Key Insight}
      Narrative is not just a product of cognition—it is an ongoing cognitive process.
    \end{block}
  }
\end{frame}

%=============================================================================
\section{Human Event Processing}
%=============================================================================

\begin{frame}{What We Know About Human Event Processing}
  % Add content about event segmentation, event cognition research, etc.
  
  \begin{itemize}
    \item Event segmentation theory
    \item Predictive processing in narrative comprehension
    \item The role of schemas and scripts
  \end{itemize}
  
  % Placeholder for future content
\end{frame}

%=============================================================================
\section{AI Event Processing: State of the Art}
%=============================================================================

\begin{frame}{Current Approaches in AI}
  % Add content about current AI approaches to event/narrative processing
  
  \begin{columns}
    \column{0.5\textwidth}
    \textbf{Narrative Analysis}
    \begin{itemize}
      \item Story understanding systems
      \item Narrative structure extraction
      \item Character/event modeling
    \end{itemize}
    
    \column{0.5\textwidth}
    \textbf{Narrative Generation}
    \begin{itemize}
      \item Large language models
      \item Story generation systems
      \item Interactive narrative
    \end{itemize}
  \end{columns}
\end{frame}

%=============================================================================
\section{Research Roadmap}
%=============================================================================

\begin{frame}{Towards Human-like Event Processing in AI}
  % Research directions
  
  \begin{enumerate}
    \item Direction 1: [To be developed]
    \item Direction 2: [To be developed]
    \item Direction 3: [To be developed]
  \end{enumerate}
  
  \vspace{1em}
  
  \textbf{Contribution to AGI:} [To be developed]
\end{frame}

%=============================================================================
\section{Ethics and Open Questions}
%=============================================================================

\begin{frame}{Is Embodiment Essential?}
  \begin{block}{Question 1}
    Is embodiment an essential precondition of narrative processing?
  \end{block}
  
  % Add discussion points
\end{frame}

\begin{frame}{Limitation or Feature?}
  \begin{block}{Question 2}
    Is the ``functionality'' of narrative processing actually a limitation imposed by processing capacity?
  \end{block}
  
  \pause
  
  \begin{block}{Corollary}
    Might artificial systems ultimately transcend that limitation?
  \end{block}
  
  % Add discussion points
\end{frame}

%=============================================================================
\section{Conclusion}
%=============================================================================

\begin{frame}{Summary}
  \begin{itemize}
    \item Ongoing narrative processing is central to human cognition
    \item Current AI approaches don't capture this process
    \item Research directions exist to bridge this gap
    \item Important ethical and philosophical questions remain
  \end{itemize}
  
  \vspace{1em}
  
  \centering
  \textbf{Thank you!}
  
  \vspace{0.5em}
  
  \textit{Questions?}
\end{frame}

%=============================================================================
% References
%=============================================================================

\begin{frame}[allowframebreaks]{References}
  \printbibliography
\end{frame}

\end{document}
